\scanpage{8}
\mspage{- 5 -}
\begin{flushright}
n\textsuperscript{o}\handwritten{354}
\end{flushright}

2) Les $\rlap{\ensuremath{\handwritten{\underline{\underline{\phantom{Z}}}}}}Z$-modules gradués de type $\Delta$, lorsque $\rlap{\ensuremath{\handwritten{\underline{\underline{\phantom{Z}}}}}}Z$ est muni de la graduation triviale, ne sont autres que les groupes commutatifs gradués de la déf. 1.

DEFINITION 4.- \uline{Soient} $A$, $A'$ \uline{deux anneaux gradués de même type} $\Delta$, $(A_\lambda)$, $(A'_\lambda)$ \uline{leurs graduations respectives}. \uline{On dit qu'un homomorphisme d'anneaux} $h:A\to A'$ \uline{est gradué si l'on a} $h(A_\lambda)\subset A'_\lambda$ \uline{pour tout} $\lambda\in\Delta$. \uline{Soient} $M$, $M'$ \uline{deux modules gradués de type} $\Delta$ \uline{sur un anneau gradué} $A$ \uline{de type} $\Delta$, $(M_\lambda)$, $(M'_\lambda)$ \uline{les graduations respectives de} $M$ \uline{et} $M'$. \uline{On dit qu'un} $A$-\uline{homomorphisme} $u:M\to M'$ \uline{est gradué de degré} $\delta\in\Delta$ \uline{si l'on a} $u(M_\lambda)\subset M'_{\lambda+\delta}$ \uline{pour tout} $\lambda\in\Delta$.

Si $u\ne 0$ et si tout élément de $\Delta$ est \uline{régulier}, le degré $\delta$ de $u$ est déterminé de façon \uline{unique}. Si $h:A\to A'$, $h':A'\to A''$ sont deux homomorphismes gradués d'anneaux gradués de type $\Delta$, il en est de même de $h'\circ h:A\to A''$ ; pour qu'une application $h:A\to A'$ soit un isomorphisme d'anneaux gradués, il faut et il suffit que $h$ soit bijective et que $h$ et l'application récipro- $h'$ soient des homomorphismes gradués ; il suffit d'ailleurs pour cela que $h$ soit un homomorphisme gradué bijectif. On voit donc qu'on peut prendre pour \uline{morphismes} de l'espèce de structure d'anneau gradué de type $\Delta$ les homomorphismes gradués (\uline{Ens}., Chap.IV, \S{} 2,n\textsuperscript{o}1).

De même, si $u:M\to M'$ et $v:M'\to M''$ sont deux homomorphismes gradués $\ne 0$ de $A$-modules gradués de type $\Delta$, de degrés respectifs $\delta$ et $\delta'$, $u'\handwritten{\unsure{\circ|\mathrm{o}}}u:M\to M''$ est un homomorphisme gradué de degré $\delta+\delta'$. Si $\delta$ est inversible et si $u:M\to M'$ est un homomorphisme gradué bijectif de degré $\delta$, l'application réciproque $u':M'\to M$ est un homomorphisme gradué bijectif de degré $-\delta$. On en conclut comme ci-dessus qu'on peut prendre pour \uline{morphismes} de l'espèce de structure de $A$-module gradué de type $\Delta$, les \uline{homomorphismes gra}-
% archive sheet number: 7

